\begin{center}\begin{tikzpicture}[font=\small,>=stealth,line join=round]
\draw[gray!40,rounded corners=3pt] (0,0) rectangle (14.2,8.8);\node at (7.1,8.200000000000001) {易错三：经过原点，与在(0,1)有空缺是两件事};
\node at (3.1500000000000004,7.3) {$y=x^2$};\begin{scope}[shift={(3.1500000000000004,3.0428571428571427)},x=1.225cm,y=1.0857142857142856cm]\draw[->,gray] (-2,0)--(2,0) node[right]{$x$};\draw[->,gray] (0,-0.5)--(0,3) node[above]{$y$};\draw[blue,thick,solid,domain=-1.5:1.5,samples=180] plot(\x,{(\x)^2});\fill[orange!85!black] (0,0) circle(2pt);\end{scope}\node[align=center,text width=5.8cm] at (3.1500000000000004,1.7) {正指数的完整图像\\包含$(0,0)$};\node at (9.95,7.3) {$y=x^0$};\begin{scope}[shift={(9.95,3.0428571428571427)},x=1.225cm,y=1.0857142857142856cm]\draw[->,gray] (-2,0)--(2,0) node[right]{$x$};\draw[->,gray] (0,-0.5)--(0,3) node[above]{$y$};\draw[gray,thick,solid,domain=-1.6:-0.02,samples=180] plot(\x,{1});\draw[gray,thick,solid,domain=0.02:1.6,samples=180] plot(\x,{1});\draw[gray,fill=white,thick] (0,1) circle(3pt);\end{scope}\node[align=center,text width=5.8cm] at (9.95,1.7) {$x\ne0$时为1\\空缺位置是$(0,1)$};
\node[text width=12.6cm,align=center] at (7.1,.65) {两图同尺度；右图的空缺在水平线$y=1$上的$(0,1)$。};
\end{tikzpicture}\end{center}
